\begin{figure}[htbp]
\centering
\resizebox{\linewidth}{!}{%
\begin{tikzpicture}[font=\scriptsize]

% ---- lado radio ----
\node[etiqueta, font=\scriptsize\bfseries\color{ambar}, anchor=west] at (0,2.35) {POR EL AIRE (LoRa): lo que emite el nodo};
\foreach \i/\b in {0/12,1/34,2/18,3/11,4/58,5/01,6/00,7/37}
  \node[campo, minimum width=6.2mm, minimum height=6mm, fill=azuloscuro!8, font=\ttfamily\scriptsize] (b\i) at (\i*0.62+0.31,1.6) {\b};
\foreach \i/\b in {8/{[},9/T,10/2,11/3,12/.,13/4,14/{]}}
  \node[campo, minimum width=6.2mm, minimum height=6mm, fill=ambar!12, draw=ambar, font=\ttfamily\scriptsize] (b\i) at (\i*0.62+0.31,1.6) {\b};
\draw[decorate, decoration={brace, mirror, amplitude=3pt}, color=azuloscuro] (b0.south west) -- (b7.south east)
  node[midway, below=4pt, font=\scriptsize\color{azuloscuro}] {evento Contact ID · 8 bytes BCD};
\draw[decorate, decoration={brace, mirror, amplitude=3pt}, color=ambar] (b8.south west) -- (b14.south east)
  node[midway, below=4pt, font=\scriptsize\color{ambar}] {valor · opcional};

% ---- plataforma ----
\node[cajaacento, text width=3.6cm, fill=verdeacento!6] (plat) at (2.5,-1.9)
  {\textbf{PLATAFORMA}\\[2pt]\scriptsize recibe vía ChirpStack + MQTT,\\ separa evento y valor,\\ los guarda en dos planos};
\draw[flechagruesa] (2.5,0.5) -- (plat.north);

% ---- lado internet ----
\node[caja, text width=5.0cm, align=left, draw=azuloscuro, line width=1pt, font=\ttfamily\scriptsize, anchor=west] (dc09) at (7.6,-1.9)
  {\textrm{\bfseries\color{azuloscuro} POR INTERNET: mensaje DC-09}\\[2pt]
   token \ \ \ : *ADM-CID\\
   msgdata : 1234181158010037\\
   extdata : [T23.4]\\
   \textrm{\scriptsize\color{gristexto} + cifrado AES, secuencia, CRC}};
\draw[flecha] (plat.east) -- node[nota, below=2pt, font=\tiny, text width=2.6cm, align=center] {solo si reenvía a una\\central de alarmas externa} (dc09.west);

\node[notaclave, text width=13.4cm, align=left, font=\footnotesize] at (6.4,-5.0)
  {\textbf{El nodo nunca genera DC-09 ni toca MQTT}: llena un arreglo de bytes y la librería LoRaWAN del chip
   se encarga de la radio. «Datos extendidos» es el nombre que DC-09 da al concepto; en el aire son
   simplemente los bytes que siguen al evento. \textbf{La plataforma es la traductora} entre ambas envolturas.
   El formato exacto de los bytes del valor en el aire se fija en la Etapa 2.};
\end{tikzpicture}}
\caption{Un mismo contenido, dos envolturas: el formato compacto en el aire y el mensaje DC-09 en
internet. La plataforma traduce entre ambos y separa los dos planos.}
\label{fig:envolturas}
\end{figure}
